\documentclass[11pt]{article}
\usepackage[margin=1in]{geometry}
\usepackage{amsmath,amssymb,amsthm,mathtools,microtype,hyperref}
\hypersetup{hidelinks}
\newtheorem{theorem}{Theorem}
\newtheorem{lemma}{Lemma}
\newcommand{\ddc}{dd^c}
\newcommand{\PP}{\mathbb P}
\newcommand{\R}{\mathbb R}
\newcommand{\V}{\mathcal V}
\title{Integer degeneration orders in nonclosed projective background cones}
\author{}
\date{4 October 2026}
\begin{document}
\maketitle
\begin{abstract}
We specify cones of smooth semipositive, possibly nonclosed backgrounds
on products of projective spaces and their linear blowups. Every
background in these fixed cones has a common integer degeneration
order for its Gauduchon variational volume, independently of the
positive smooth probability and perturbation metric. Quantitative
bounds follow from the full Gauduchon infimum: an explicit logarithmic
exceptional-divisor potential compares the background with a closed
constant-rank form. Product linear blowups preserve the order. For
specified effective real divisors satisfying explicit cone-preserving
parameter conditions, we give a necessary and sufficient equality criterion and the
exact integer gap. Nonclosed exceptional-divisor examples realize
every jump from one to the dimension minus one. An extension to
arbitrary modifications with K\"ahler source uses the known closed-nef
order theorem. These are sufficient conditions for P02, rather than
a classification of arbitrary nonclosed backgrounds.
\end{abstract}

\section{The volume and a full-infimum comparison}
Throughout use $\ddc=(i/(2\pi))\partial\bar\partial$, so
$\ddc\log|z|^2=\delta_0$ on a transverse disk. Let $M$ be a
connected compact complex manifold of dimension $n$, $\mu$ a smooth
strictly positive probability volume and $h$ a smooth positive Hermitian
form. For a smooth real background $A$ admitting a positive smooth
potential representative, put
\[
 D_\mu(\gamma)=\int_M
       \left(\frac{\mu}{\gamma^n}\right)^{1/n}\gamma^n,
 \qquad
 \lambda_\mu(A)=\inf_{\gamma}
       \frac{\int_M A\wedge\gamma^{n-1}}{D_\mu(\gamma)},
 \qquad \V_\mu(A)=\lambda_\mu(A)^n.
\]
The infimum ranges over \emph{all} smooth positive Hermitian forms
$\gamma$ with $\ddc\gamma^{n-1}=0$. No invariant or product
restriction is imposed on these tests. When $n=1$ the denominator is
one, and $\lambda_\mu(A)=\int_M A$.

A background $\theta$ is nef here if for every $\varepsilon>0$,
$\theta+\varepsilon h$ has a positive smooth potential representative.
Its lower and upper degeneration orders are the liminf and limsup of
\[
 \frac{\log\V_\mu(\theta+\varepsilon h)}{\log\varepsilon}
 \qquad(\varepsilon\downarrow0).
\]
The quantitative results below prove that these two orders coincide.
They concern this variational volume, rather than extrema of
Monge--Amp\`ere masses over bounded potentials.

\begin{lemma}\label{lem:sandwich}
Suppose $M$ has a K\"ahler form $\Omega$, and $B$ is a smooth
closed semipositive $(1,1)$-form. Let $\theta$ be smooth nef and
assume, for constants $c,C>0$ and a global real $\varphi\in L^1(M)$,
\[
 \theta\leq CB,\qquad \theta+\ddc\varphi\geq cB.
\]
There are constants $k,K>0$, independent of $0<\varepsilon\leq1$,
such that
\[
 k\lambda_\mu(B+\varepsilon\Omega)
 \leq\lambda_\mu(\theta+\varepsilon h)
 \leq K\lambda_\mu(B+\varepsilon\Omega).
\]
If $B$ has constant complex rank $r$, there are $L,U>0$ with
\[
 L\varepsilon^{n-r}\leq\V_\mu(\theta+\varepsilon h)
                  \leq U\varepsilon^{n-r}.
\]
In particular both orders equal $n-r$.
\end{lemma}
\begin{proof}
Choose $c_h,C_h>0$ with $c_h\Omega\leq h\leq C_h\Omega$.
For every smooth Gauduchon test, distributional integration by parts
gives
\[
 \int_M\ddc\varphi\wedge\gamma^{n-1}
       =\int_M\varphi\ddc\gamma^{n-1}=0.
\]
The integrals are well defined because $\varphi$ is locally integrable
and the test is smooth. Positivity of the current lower bound implies
\[
 \int_M(\theta+\varepsilon h)\wedge\gamma^{n-1}
 \geq\min(c,c_h)\int_M(B+\varepsilon\Omega)
                                      \wedge\gamma^{n-1}.
\]
The pointwise upper bound gives the reverse inequality with
$\max(C,C_h)$. Divide by the identical positive denominator and
take both infima over the entire test cone.

For the constant-rank assertion set $A_\varepsilon=B+\varepsilon\Omega$.
The $r$ positive eigenvalues of $B$ relative to $\Omega$ have a
uniform positive lower bound and finite upper bound. Compactness and
the positive smooth density $\mu$ give constants $l,u>0$ such that
\begin{equation}\label{eq:pointwise-volume}
 l\varepsilon^{n-r}\mu\leq A_\varepsilon^n
                         \leq u\varepsilon^{n-r}\mu.
\end{equation}
For any positive test, matrix arithmetic--geometric mean gives
\[
 A_\varepsilon\wedge\gamma^{n-1}
 \geq\left(\frac{A_\varepsilon^n}{\gamma^n}\right)^{1/n}
                                                      \gamma^n.
\]
The lower bound in \eqref{eq:pointwise-volume}, integrated against
every Gauduchon test, proves
$\lambda_\mu(A_\varepsilon)\geq
 l^{1/n}\varepsilon^{(n-r)/n}$.
For the upper bound use the admissible test $\gamma=A_\varepsilon$,
which is positive and closed. If $A_\varepsilon^n=f_\varepsilon\mu$,
then
\[
 D_\mu(A_\varepsilon)=\int_Mf_\varepsilon^{(n-1)/n}\mu
 \geq l^{(n-1)/n}\varepsilon^{(n-r)(n-1)/n},
 \qquad \int_M A_\varepsilon^n\leq u\varepsilon^{n-r}.
\]
Thus $\lambda_\mu(A_\varepsilon)\leq
 ul^{-(n-1)/n}\varepsilon^{(n-r)/n}$. Raising to the $n$th power
and applying the comparison proves the quantitative assertion,
including $r=0$ and $n=1$. Fixed multiplicative constants disappear
after taking logarithms and dividing by $\log\varepsilon$.
\end{proof}

\section{The geometry and the fixed cone}
A plain factor is $Z=\PP^m$, $m\geq1$, with normalized
Fubini--Study form $H$ and $\kappa=H$. A blown factor is
\[
 Z=\operatorname{Bl}_{\PP^\ell}\PP^m,
 \qquad m\geq2,\quad0\leq\ell\leq m-2,
 \qquad r=m-\ell-1.
\]
Writing projective coordinates as $[z_0:\cdots:z_\ell:w_0:\cdots:w_r]$,
its explicit model is
\[
 Z=\{(x,[t])\in\PP^m\times\PP^r:
                        w_at_b=w_bt_a\text{ for all }a,b\}.
\]
Let $p:Z\to\PP^m$ and $q:Z\to\PP^r$ be the projections. In a
chart $t_j=1$ the equations are $w_a=t_aw_j$, the blowup charts
of the ideal $(w_0,\ldots,w_r)$. Equivalently write $w=\lambda t$;
the fiber coordinates $[z_0:\cdots:z_\ell:\lambda]$ exhibit $q$
as a locally trivial $\PP^{\ell+1}$-bundle. These descriptions
prove smoothness and that $q$ is a submersion.

Set $H=p^*\omega_m$, $Q=q^*\omega_r$, and $\kappa=H+Q$.
Both $H,Q$ are closed semipositive, and $Q$ has constant rank $r$.
The incidence inclusion is a smooth holomorphic embedding in
$\PP^m\times\PP^r$, so its restricted product form $\kappa$
is K\"ahler. The exceptional divisor $E=\{w=0\}$ satisfies
\[
 \mathcal O(E)=p^*\mathcal O(1)\otimes q^*\mathcal O(-1).
\]
Indeed the coordinate section $w_j$ has divisor
$E+q^{-1}\{t_j=0\}$, as is also seen from $w_j=\lambda t_j$.
Use the indicated product metric of curvature $H-Q$. For the
canonical section $s_E$, put $\psi_E=\log|s_E|^2$. The transverse
disk normalization and the smooth metric term give the
Poincar\'e--Lelong identity \cite[III.2.15 and V.13.2]{Demailly}
\begin{equation}\label{eq:exceptional}
 \ddc\psi_E=[E]-H+Q,
 \qquad H+\ddc(\psi_E/2)=\tfrac12(\kappa+[E])\geq\kappa/2.
\end{equation}
These logarithms are locally integrable. The identity supplies the
current lower bound without an analytic existence theorem.

Let $X=\prod_{i=1}^s Z_i$ be a product of these factors,
$m_i=\dim Z_i$, $n=\sum_i m_i$, with all factor forms understood
as pullbacks. Fix once and for all smooth strictly positive functions
$f_i$ on $X$, and $g_i$ for the blown factors. They may depend on
every product coordinate. Define the fixed cone
\begin{equation}\label{eq:cone}
 \theta=\sum_i\bigl((a_i+b_if_i)H_i+(c_i+d_ig_i)Q_i\bigr),
 \qquad a_i,b_i,c_i,d_i\geq0,
\end{equation}
omitting the $Q_i$ terms and their parameters on plain factors.
All functions and generators are fixed independently of the parameters,
the perturbation, the probability and the divisors.

Every background is smooth semipositive. Thus $\theta+\varepsilon h$
itself is positive for every $h>0$ and $\varepsilon>0$; the precise
smooth-nef condition is verified. Say $H_i$ is active if $a_i+b_i>0$,
and $Q_i$ is active if $c_i+d_i>0$. Strict positivity of the fixed
functions gives a uniform positive lower bound for every active
coefficient. Define
\[
 R_i(\theta)=
 \begin{cases}
 m_i,&H_i\text{ active},\\
 r_i=m_i-\ell_i-1,&H_i\text{ inactive and }Q_i\text{ active},\\
 0,&\text{both inactive},
 \end{cases}
 \qquad R(\theta)=\sum_iR_i(\theta),
\]
where the middle alternative applies only to blown factors.

\begin{theorem}\label{thm:orders}
For every background in \eqref{eq:cone}, every positive smooth
$h$ and every strictly positive smooth probability $\mu$, there
are constants $L,U>0$ such that
\[
 L\varepsilon^{n-R(\theta)}
 \leq\V_\mu(\theta+\varepsilon h)
 \leq U\varepsilon^{n-R(\theta)}\qquad(0<\varepsilon\leq1).
\]
Its common degeneration order is the integer $n-R(\theta)$.
\end{theorem}
\begin{proof}
Let $B_\theta$ be the sum of $\kappa_i$ on $H_i$-active blocks,
$Q_i$ on $Q_i$-only blocks, and zero on inactive blocks. It is
closed semipositive of constant rank $R(\theta)$. The form
$\Omega=\sum_i\kappa_i$ is K\"ahler. Boundedness of the fixed
coefficients and $H_i,Q_i\leq\kappa_i$ imply
$\theta\leq CB_\theta$.

On each $H_i$-active blown block put
$u_i=\min_X(a_i+b_if_i)>0$, and set
\[
 \varphi=\sum_{\substack{i\text{ blown}\\H_i\text{ active}}}
                                                   (u_i/2)\psi_{E_i}.
\]
By \eqref{eq:exceptional}, the corresponding $u_iH_i$ becomes
at least $(u_i/2)\kappa_i$, with a positive divisor current. The
remainder $(a_i+b_if_i-u_i)H_i$ and all $Q_i$ terms remain
nonnegative. On plain active blocks use the positive minimum of
$a_i+b_if_i$; on $Q_i$-only blocks use that of $c_i+d_ig_i$.
The minimum of these finitely many constants gives
$\theta+\ddc\varphi\geq cB_\theta$.
For the zero background use $\varphi=0$. Lemma~\ref{lem:sandwich}
proves the assertion from the full variational definition.
\end{proof}

The exterior derivative of a cone member is
\[
 d\theta=\sum_i b_i\,df_i\wedge H_i
       +\sum_{i\text{ blown}}d_i\,dg_i\wedge Q_i.
\]
If a fixed generator derivative is nonzero, activating that generator
alone gives a genuinely nonclosed member. Such choices of the fixed
functions exist. On a blown factor choose $g>0$ with a nonzero
derivative along a $q$-fiber, using a small
smooth bump perturbation of one in a local product chart. Then
$dg\wedge Q\ne0$, so $gQ$ is nonclosed and has order $\ell+1$.
On $\PP^p\times\PP^q$, $p,q\geq1$, take $f>0$ varying
nontrivially on the second factor. The background $fH_1$ is
nonclosed and has order $q$. These exterior derivatives cannot
be removed by adding $\ddc$ of a smooth function, because
$d\ddc=0$. The manifold being projective does not make its
background closed.

\section{Modification equality}
Replace any collection of plain $\PP^{m_i}$ factors by the indicated
linear blowups, leaving all other factors unchanged. Let $\pi:Y\to X$
be this product modification. Fix the old coefficient functions by
pullback; choose any positive functions for the new $Q$ generators.
The pulled-back old cone lies in the new fixed cone with all new
$Q$ coefficients zero. An active plain $H_i$ becomes an active blown
$H_i$ of the same effective rank $m_i$; inactive blocks remain
inactive. Theorem~\ref{thm:orders} on both products gives
\begin{equation}\label{eq:modification}
 a_-(\pi^*\theta)=a_+(\pi^*\theta)
                     =a_-(\theta)=a_+(\theta)=n-R(\theta).
\end{equation}
Every allowed modification satisfies equality, including when the
pointwise rank of a pulled-back $H_i$ drops on its exceptional divisor.
The logarithmic-current lower bound accounts for that drop. Both
sides have the proved quantitative power bounds.

There is a further extension with a stated additional premise.
Use the known theorem that on a compact K\"ahler manifold the
degeneration orders of a smooth closed nef form equal the dimension
minus its numerical dimension. Then \eqref{eq:modification} holds
for every holomorphic modification $\pi:Y\to X$ with $Y$ compact
K\"ahler. To see the entire reduction, pull back the explicit
logarithmic identities and nonnegative smooth remainders in the proof
above. Nonzero divisor sections remain nonzero holomorphic sections
under the surjective modification; their logarithms remain locally
integrable. Hence
\[
 \pi^*\theta\leq C\pi^*B_\theta,
 \qquad \pi^*\theta+\ddc(\varphi\circ\pi)\geq c\pi^*B_\theta.
\]
The comparison part of Lemma~\ref{lem:sandwich}, which does not
require constant rank, compares the full infimum with that of
$\pi^*B_\theta+\varepsilon\Omega_Y$ for a K\"ahler form
$\Omega_Y$. Put $R=R(\theta)$. Then
$(\pi^*B_\theta)^{R+1}=0$ identically, whereas
\[
 \int_Y(\pi^*B_\theta)^R\wedge\Omega_Y^{n-R}>0.
\]
The latter is positive over a nonempty open set where $\pi$ is an
isomorphism; for $R=0$ it is the positive K\"ahler volume.
Thus its numerical dimension is $R$. The stated closed-nef theorem
proves the extension. This step claims common orders, and does not
assert quantitative two-sided power bounds for arbitrary $Y$ from
that premise alone. The product-modification result needs no such
external order theorem. Arbitrary non-K\"ahler sources are excluded.

\section{Effective divisors and the exact equality criterion}
Fix the line bundles $\mathcal O(H_i)$ and, on blown factors,
$\mathcal O(Q_i)$ with the pullback Fubini--Study metrics used above.
Consider tensor products of their integer powers, permitting negative
powers, with these induced metrics. Let
\[
 D=\sum_\nu t_\nu\operatorname{div}(s_\nu),\qquad t_\nu\geq0,
\]
where the sum is finite and every $s_\nu$ is a nonzero holomorphic
section of one such tensor product. This specifies effective real divisors and their
curvature representatives:
\[
 \eta_D=\sum_i(\delta_iH_i+\rho_iQ_i),
 \qquad \Phi_D=\sum_\nu t_\nu\log|s_\nu|^2,
 \qquad [D]=\eta_D+\ddc\Phi_D\geq0.
\]
Again omit $Q$ terms on plain factors. The section norms fix the
smooth curvature; arbitrary changes of representatives are not hidden
in the definition.

The coefficient range of effective divisors is
\begin{equation}\label{eq:effective-range}
 \delta_i\geq0\quad\text{on plain factors},\qquad
 \delta_i\geq0,\quad\delta_i+\rho_i\geq0
                                      \quad\text{on blown factors}.
\end{equation}
For a nonzero section choose a point where it is nonzero and which
lies outside all exceptional divisors. Hold the other factor
coordinates fixed. A line through this point in a $q_i$-fiber
$\PP^{\ell_i+1}$ has $H_i$-degree one and $Q_i$-degree zero,
so the $H_i$ exponent is nonnegative. For the other inequality write
the chosen blown-factor point as $([z:w],[w])$, $w\ne0$. Choose
a linear map $A$ with $Aw=z$. The section of $q_i$ given by
\[
 [t]\longmapsto([At:t],[t])
\]
contains this point, and both $H_i,Q_i$ have degree one on its
projective lines. Since $r_i\geq1$, choose such a line through
the point. Its section restriction is nonzero, so the sum of the
exponents is nonnegative. The plain-factor line argument gives
the other inequality. These degree inequalities also follow directly
by integrating the section's Poincar\'e--Lelong formula on the line.
Positive real sums preserve the conditions. Conversely every pair
in the blown-factor range
is represented by
$\delta_iE_i+(\delta_i+\rho_i)$ times a $Q_i$-hyperplane divisor.
The plain range is represented by hyperplanes. Thus these are exact
numerical ranges for the specified real divisor family.

\begin{theorem}\label{thm:residual}
Let $\theta$ have the parameters in \eqref{eq:cone}, and let $D$
be as above. Suppose
\[
 a_i\geq\delta_i\text{ for every }i,
 \qquad c_i\geq\rho_i\text{ for every blown }i.
\]
Then $\beta=\theta-\eta_D$ belongs to the same fixed cone and is
smooth semipositive and nef. Its common order satisfies
\[
 a(\beta)-a(\theta)=\sum_i\bigl(R_i(\theta)-R_i(\beta)\bigr).
\]
Equality in both residual-order comparisons holds if and only if
\begin{align*}
 a_i-\delta_i+b_i&>0&&\text{whenever }a_i+b_i>0,\\
 c_i-\rho_i+d_i&>0&&\text{whenever }a_i=b_i=0,
                                  \ c_i+d_i>0.
\end{align*}
The second condition applies only on blown factors. Here $a$ denotes
the common order, and every summand in the gap is nonnegative.
\end{theorem}
\begin{proof}
The residual parameters are
$(a_i-\delta_i,b_i,c_i-\rho_i,d_i)$, all nonnegative.
Thus $\beta+\varepsilon h$ itself is positive for every $h>0$.
Moreover $(\theta-\beta)+\ddc\Phi_D=[D]\geq0$, the required
actual positive-current comparison. Theorem~\ref{thm:orders} gives
the displayed difference of orders.

If an initially active $H_i$ is erased, its effective rank drops
from $m_i$ to either $r_i<m_i$ or zero. If $H_i$ was absent, then
$a_i=b_i=0$, and \eqref{eq:effective-range} and the residual
conditions force $\delta_i=0$ and $\rho_i\geq0$. Such a block
cannot gain $Q_i$; a $Q_i$-only block can keep or lose its rank.
Therefore ranks cannot increase in any block, and rank losses cannot
cancel between blocks. Equality holds exactly when no initially
leading active generator is erased, as stated.
\end{proof}

The stated conditions ensure that the displayed residual parameters
remain nonnegative. They are sufficient conditions for this theorem,
and do not characterize all nef residual backgrounds or all alternative
representations in the cone.

\section{Nonclosed examples with sharp integer gaps}
On $Z=\operatorname{Bl}_{\PP^\ell}\PP^m$ choose fixed $g>0$
with $dg\wedge Q\ne0$. For $t>0$ put
\[
 \theta_t=tH+gQ.
\]
It is genuinely nonclosed and strictly positive, since
$\theta_t\geq\min(t,\min g)(H+Q)>0$. Its order is zero.
For $0\leq s\leq t$ remove $sE$ with the fixed curvature
$s(H-Q)$. The residual and actual boundary current are
\[
 \beta_s=(t-s)H+(g+s)Q,
 \qquad \theta_t+\ddc(s\psi_E)=s[E]+\beta_s.
\]
Every $\beta_s$ is smooth semipositive, nef and nonclosed, and
\[
 a(\beta_s)=0\quad(s<t),\qquad
 a(\beta_t)=m-r=\ell+1.
\]
Equality thus holds precisely for $s<t$. Varying $\ell$ realizes
all gaps $1,\ldots,m-1$. At $s=t$ the chosen current contains
the actual divisorial mass $t[E]$ before removal. This does not
assert a positive minimal perturbative divisorial multiplicity:
the smooth positive representatives $\theta_t+\varepsilon h$ and
$\beta_s+\varepsilon h$ show those minimal multiplicities are zero.

For a nonclosed degenerate example with modification equality take
$\theta=fH_1$ on $\PP^p\times\PP^q$, where $f>0$ varies on
the second factor. Its order is $q$, and every allowed product
linear blowup preserves it, including blowups in active or inactive
directions. The stated extension gives the same conclusion for every
modification with K\"ahler source.

\section{Scope}
These sufficient conditions and equality criteria use fixed strictly
positive coefficient functions and explicitly cone-preserving divisors
with fixed curvature metrics. The probability and perturbation metric
are arbitrary smooth positive data; every bound uses the full
Gauduchon infimum. The product results use elementary geometry and
current integration. The arbitrary-K\"ahler-modification extension
uses the closed-nef order theorem supplied in P02. Coefficient functions
with zeros, arbitrary residual cones and non-K\"ahler sources are not
covered. The general P02 programme remains open.
\begin{thebibliography}{9}
\bibitem{Demailly}
J.-P. Demailly, \emph{Complex Analytic and Differential Geometry},
version of 21 June 2012, III.2.10--2.15 and V.13.1--13.8.
\end{thebibliography}
\end{document}
